%% toc: P2 20 20.1,20.2,20.3,20.4
\section{Proofs on choice-free completeness and compactness}
\label{sec:annex-p02-w02c-p02-023-choice-free}

Notation and the ambient theory \(\mathsf{ZF}\) are those of
\cref{sec:p02-w02c-p02-023-choice-free} and
\cref{def:p02-w02c-p02-023-zf}.  ``Finite-set lemma'' refers to
\cref{thm:p01-w02-p01-011-finite-sets} and ``pigeonhole'' to
\cref{thm:p01-w02-p01-008-pigeonhole}.  Fixing one element of a set already
shown to be nonempty is a single existential instantiation; it is not an
application of the Axiom of Choice, and it is the only kind of selection the
proofs below make.

\begin{annexitem}{annex:p02-w02c-p02-023-compact}{thm:p02-w02c-p02-023-countable-compactness}{Countable compactness without choice}
\begin{proof}
\emph{Clause 1, left to right.}  If \(\mathcal M\entails T\) and
\(T_0\subseteq T\), then \(\mathcal M\entails\sigma\) for every
\(\sigma\in T_0\), so \(\mathcal M\entails T_0\).

\emph{Clause 1, right to left.}  Suppose every finite subset of \(T\) has a
model.

\emph{Step 1: \(T\) is consistent.}  Suppose an \(\mathsf{ND}_{\mathsf{FOL}}\)-derivation
\(D\) of \(\bottom\) has all its open assumptions in \(T\).  \(D\) is a
finite tree, so the set \(T_0\) of its open assumptions is a finite subset
of \(T\), and \(D\) witnesses \(T_0\derives_{\mathsf{ND}_{\mathsf{FOL}}}\bottom\).
\Cref{thm:fol-soundness} gives \(T_0\entails\bottom\).  By hypothesis some
\(\mathcal M\) is a model of \(T_0\).  Its domain is nonempty; fix
\(a\in\dom\mathcal M\) and let \(s\) be the assignment with value \(a\) at every
variable.  Each \(\gamma\in T_0\) is a sentence with
\(\mathcal M\entails\gamma\), so \(\mathcal M,s\entails\gamma\) by the definition of
truth of a sentence in \cref{sec:fol-semantics-02}.  Then
\(T_0\entails\bottom\) gives \(\mathcal M,s\entails\bottom\), which no structure
and assignment satisfy.  Hence no such \(D\) exists:
\(T\notderives_{\mathsf{ND}_{\mathsf{FOL}}}\bottom\).

\emph{Step 2: \(T\) is satisfiable.}  Suppose that no \(\Lang\)-structure
\(\mathcal M\) and assignment \(s\) satisfy every member of \(T\).  Then the
condition defining \(T\entails\bottom\) holds, because its antecedent is
never met.  The signature \(\Lang\) is countable and
\(T\cup\{\bottom\}\subseteq\operatorname{Fm}(\Lang)\), so
\cref{cor:fol-completeness-countable} gives
\(T\derives_{\mathsf{ND}_{\mathsf{FOL}}}\bottom\) with no use of the Axiom of Choice.  This
contradicts Step~1.  Therefore some \(\Lang\)-structure \(\mathcal M\) and
assignment \(s\) satisfy every member of \(T\).

\emph{Step 3: \(\mathcal M\) is a model of \(T\).}  Let \(\sigma\in T\).  Since
\(\sigma\) is a sentence, \(\FV(\sigma)=\varnothing\), so every assignment
\(s'\) agrees with \(s\) on \(\FV(\sigma)\), and \cref{lem:fol-coincidence}
gives \(\mathcal M,s'\entails\sigma\).  Thus \(\mathcal M\entails\sigma\) for every
\(\sigma\in T\), that is, \(\mathcal M\entails T\).

\emph{Clause 2.}  Suppose \(T\entails\sigma\).  Since \(\Lang\) is countable
and \(T\cup\{\sigma\}\subseteq\operatorname{Fm}(\Lang)\),
\cref{cor:fol-completeness-countable} gives an \(\mathsf{ND}_{\mathsf{FOL}}\)-derivation
\(D\) of \(\sigma\) with open assumptions in \(T\).  The set \(T_0\) of its
open assumptions is a finite subset of \(T\), and \(D\) witnesses
\(T_0\derives_{\mathsf{ND}_{\mathsf{FOL}}}\sigma\).  \Cref{thm:fol-soundness} gives
\(T_0\entails\sigma\).

\emph{Choice.}  The proof uses \cref{cor:fol-completeness-countable}, whose
statement records that it makes no use of the Axiom of Choice; soundness,
which is an induction on derivations; \cref{lem:fol-coincidence}, a
structural induction; and the single element \(a\) fixed in Step~1.  No
family of selections is made.
\end{proof}
\end{annexitem}

\begin{annexitem}{annex:p02-w02c-p02-023-finite}{thm:p02-w02c-p02-023-finite-structures}{Finiteness unaxiomatizable without choice}
\begin{proof}
In the empty signature a structure is a nonempty domain and nothing else.
For \(n\geq2\) let \(\iota_n\) and \(T_\infty=\{\iota_n:n\geq2\}\) be as in
the example on infinite structures in \cref{ch:fol-theories}.  For
\(N\geq2\) let \(\mathcal E_N\) be the structure with domain
\(N=\{0,\ldots,N-1\}\).

\emph{Claim A: \(\mathcal E_N\entails\iota_n\) whenever
\(2\leq n\leq N\).}  Let \(s\) be any assignment in \(\mathcal E_N\) and put
\(s'=s[x_0\mapsto0]\cdots[x_{n-1}\mapsto n-1]\); the values \(0,\ldots,n-1\)
lie in \(N\) because \(n\leq N\).  The variables \(x_0,\ldots,x_{n-1}\) are
distinct members of \(\operatorname{Var}\), so \(s'(x_i)=i\) for \(i<n\).  For
\(0\leq i<j<n\) we have \(s'(x_i)\ne s'(x_j)\), so \(s'\) satisfies each
conjunct \(x_i\ne x_j\) and hence the finite conjunction.  Applying the
existential clause of satisfaction \(n\) times, from the innermost
quantifier outwards, gives \(\mathcal E_N,s\entails\iota_n\).

\emph{Claim B: if \(\dom\mathcal M\) is finite, then some \(\iota_n\) is false in
\(\mathcal M\).}  Let \(g:k\to\dom\mathcal M\) be a bijection with \(k\in\N\), and put
\(n=k+2\geq2\).  If \(\mathcal M\entails\iota_n\), the existential clause
applied \(n\) times gives an assignment \(s\) satisfying every
\(x_i\ne x_j\) with \(i<j<n\), so the map \(i\mapsto s(x_i)\) is an
injection \(n\to\dom\mathcal M\).  Composing with \(g^{-1}\) gives an injection
\(n\to k\) with \(k<n\), which pigeonhole, item~2, excludes.

\emph{The contradiction.}  Suppose \(T\) is a theory in the empty signature
whose models are exactly the structures with finite domain.  Let
\(F\subseteq T\cup T_\infty\) be finite.  The map \(n\mapsto\iota_n\) is
injective, so \(\{n:\iota_n\in F\}\) is in bijection with a subset of \(F\)
and is finite by the finite-set lemma, item~3.  A finite set of natural
numbers has a largest member when it is nonempty, by induction on its size
using item~1 of the same lemma.  Let \(N\) be the largest member of
\(\{2\}\cup\{n:\iota_n\in F\}\).  The domain \(N\) of \(\mathcal E_N\) is
finite, so \(\mathcal E_N\entails T\), and by Claim A it satisfies every
\(\iota_n\in F\).  Hence \(\mathcal E_N\entails F\): every finite subset of
\(T\cup T_\infty\) has a model.

The empty signature is countable, since the empty function is an injection
of its empty symbol set into \(\N\).  Clause~1 of
\cref{thm:p02-w02c-p02-023-countable-compactness} gives a model \(\mathcal M\)
of \(T\cup T_\infty\), with no use of the Axiom of Choice.  As
\(\mathcal M\entails T\), its domain is finite, and Claim B gives some
\(\iota_n\in T_\infty\) false in \(\mathcal M\), contradicting
\(\mathcal M\entails T_\infty\).  So no such \(T\) exists.
\end{proof}
\end{annexitem}

\begin{annexitem}{annex:p02-w02c-p02-023-bpi}{thm:p02-w02c-p02-023-bpi}{Compactness for all signatures gives BPI}
\begin{proof}
Let \(V\) be a set and \(\Gamma\) a set of propositional formulas over the
variables \(p_v\) (\(v\in V\)) in the sense of
\cref{def:p01-bpi-propositional-compactness}, such that every finite subset
of \(\Gamma\) is satisfied by some valuation \(\nu:V\to\mathbf2\).  We find
one valuation satisfying all of \(\Gamma\).

\emph{The signature.}  Let \(W\) be the union of \(\operatorname{Var}\) and the set of
reserved logical and punctuation tokens of
\cref{sec:fol-syntax-semantics-01}, and for \(v\in V\) put
\(R_v=\{\{v\},\{v,W\}\}\), the Kuratowski pair of \(v\) and \(W\).
Distinct \(v\) give distinct \(R_v\), by the characteristic property of the
pair.  No \(R_v\) lies in \(W\): otherwise \(W\in\{v,W\}\in R_v\in W\), and
the set \(\{W,\{v,W\},R_v\}\) would have no \(\in\)-minimal member,
contrary to the Axiom of Foundation.  Let \(\Lang_V\) be the signature with
no function symbols, relation symbols \(\{R_v:v\in V\}\) (a set by
Replacement), and arity zero for every \(R_v\).  An \(\Lang_V\)-structure
\(\mathcal M\) assigns each \(R_v\) a subset of \((\dom\mathcal M)^0\), whose only
member is the empty tuple \(\langle\rangle\); and \(\mathcal M,s\entails R_v\) iff
\(\langle\rangle\in R_v^{\mathcal M}\).

\emph{The translation.}  Define \(\varphi^{*}\) by recursion on
propositional formulas: \(p_v^{*}=R_v\), \(\bot^{*}=\bottom\),
\(\top^{*}=(\bottom\to\bottom)\), \((\neg\varphi)^{*}=\neg\varphi^{*}\),
\((\varphi\land\psi)^{*}=(\varphi^{*}\land\psi^{*})\) and
\((\varphi\lor\psi)^{*}=(\varphi^{*}\lor\psi^{*})\).  Each
\(\varphi^{*}\) contains no variable, so it is an \(\Lang_V\)-sentence.  The
map \(\varphi\mapsto\varphi^{*}\) is injective, by induction on
\(\varphi\): by unique readability of first-order formulas
(\cref{sec:fol-syntax-semantics-02}), \(\varphi^{*}\) is exactly one of an
atomic formula \(R_v\), the formula \(\bottom\), a negation, or a binary
formula with main connective \(\to\), \(\land\) or \(\lor\), and these cases
occur exactly when \(\varphi\) is \(p_v\), \(\bot\), a negation, \(\top\), a
conjunction or a disjunction respectively; \(R_v\) determines \(v\); and
unique readability recovers the immediate subformulas, to which the
induction hypothesis applies.

\emph{Structures and valuations.}  For an \(\Lang_V\)-structure \(\mathcal M\)
put \(\nu_{\mathcal M}(v)=1\) if \(\langle\rangle\in R_v^{\mathcal M}\) and
\(\nu_{\mathcal M}(v)=0\) otherwise.  For every propositional formula
\(\varphi\) and every assignment \(s\),
\[
 \mathcal M,s\entails\varphi^{*}
 \quad\text{iff}\quad
 \nu_{\mathcal M}(\varphi)=1,
\]
where \(\nu_{\mathcal M}\) is extended by the Boolean truth tables.  The proof is
by induction on \(\varphi\): for \(p_v\) both sides say
\(\langle\rangle\in R_v^{\mathcal M}\); \(\bottom\) is satisfied by no structure and
\(\bot\) has value \(0\); \((\bottom\to\bottom)\) is satisfied everywhere
and \(\top\) has value \(1\); and the satisfaction clauses for \(\neg\),
\(\land\) and \(\lor\) match their truth tables.  Conversely, for a
valuation \(\nu\) let \(\mathcal M_\nu\) be the \(\Lang_V\)-structure with domain
\(\{0\}\) and \(R_v^{\mathcal M_\nu}=\{\langle\rangle\}\) if \(\nu(v)=1\),
\(R_v^{\mathcal M_\nu}=\varnothing\) if \(\nu(v)=0\).  It is defined by
Separation, and \(\nu_{\mathcal M_\nu}=\nu\).

\emph{Finite satisfiability of the translation.}  Put
\(\Gamma^{*}=\{\varphi^{*}:\varphi\in\Gamma\}\subseteq\operatorname{Sent}(\Lang_V)\).  Let
\(F\subseteq\Gamma^{*}\) be finite and
\(\Gamma_F=\{\varphi\in\Gamma:\varphi^{*}\in F\}\).  By injectivity of the
translation, \(\varphi\mapsto\varphi^{*}\) is a bijection from \(\Gamma_F\)
onto \(F\), so \(\Gamma_F\) is finite.  By hypothesis some valuation \(\nu\)
satisfies every member of \(\Gamma_F\).  By the displayed equivalence,
applied to \(\mathcal M_\nu\) with \(\nu_{\mathcal M_\nu}=\nu\), the structure
\(\mathcal M_\nu\) satisfies \(\varphi^{*}\) under every assignment for each
\(\varphi\in\Gamma_F\), so \(\mathcal M_\nu\entails F\).

\emph{Conclusion.}  The assumed compactness for \(\Lang_V\) gives a model
\(\mathcal M\) of \(\Gamma^{*}\).  Fix \(a\in\dom\mathcal M\) and the assignment \(s\)
with constant value \(a\).  For \(\varphi\in\Gamma\) we have
\(\mathcal M,s\entails\varphi^{*}\), so \(\nu_{\mathcal M}(\varphi)=1\) by the
displayed equivalence.  Thus the valuation \(\nu_{\mathcal M}\), defined by
Separation from \(\mathcal M\), satisfies all of \(\Gamma\).  This is the
principle of \cref{def:p01-bpi-propositional-compactness} for \(V\), and
\cref{thm:p01-bpi-propositional-compactness} gives the forms of the Boolean
Prime Ideal Theorem in \cref{thm:p01-bpi-equivalence}, over
\(\mathsf{ZF}\).  When \(V=\varnothing\), \(\Lang_V\) is the empty signature
and the same argument applies.
\end{proof}
\end{annexitem}
